\documentclass{article}
\usepackage{listings}
\usepackage{xcolor}
\usepackage{textcomp} % For \textquotesingle

% Define colors for SQL code
\definecolor{codegreen}{rgb}{0,0.6,0}
\definecolor{codegray}{rgb}{0.5,0.5,0.5}
\definecolor{codepurple}{rgb}{0.58,0,0.82}
\definecolor{backcolour}{rgb}{0.95,0.95,0.92}
 
% Configure the listings package for SQL syntax highlighting
\lstdefinestyle{mystyle}{
    backgroundcolor=\color{backcolour},   
    commentstyle=\color{codegreen},
    keywordstyle=\color{blue},
    stringstyle=\color{codepurple},
    basicstyle=\ttfamily\footnotesize,
    breakatwhitespace=false,         
    breaklines=true,                 
    captionpos=b,                    
    keepspaces=true,                 
    numbers=left,                    
    numbersep=5pt,                  
    showspaces=false,                
    showstringspaces=false,
    showtabs=false,                  
    tabsize=2
}

\lstset{style=mystyle}

\title{SQLite Tutorial: Managing a Library System}
\author{Dr. Frank Xu}
\date{\today}

\begin{document}

\maketitle

\section*{Introduction}

In this tutorial, we will create a simple SQLite database to manage a library system. The database will store information about books, including their title, author, publication year, and genre.

\section*{Step 1: Create a New SQLite Database}

Open your terminal or command prompt and create a new SQLite database called \texttt{library.db}:

\begin{lstlisting}[language=bash]
sqlite3 library.db
\end{lstlisting}

Once you run the above command, you should see something like this in your terminal:

\begin{lstlisting}[language=bash]
user@linux:~$ sqlite3 library.db
SQLite version 3.37.2 2022-01-06 13:25:41
Enter ".help" for usage hints.
sqlite> 
\end{lstlisting}

This indicates that SQLite has successfully created the \texttt{library.db} file and opened the SQLite shell. You can now start executing SQL commands directly in this shell.

\section*{Step 2: Create a Table}

Now, let's create a table called \texttt{books} to store information about the books in our library. The table will have the following columns:

\begin{itemize}
    \item \texttt{id}: A unique identifier for each book (Primary Key).
    \item \texttt{title}: The title of the book.
    \item \texttt{author}: The author of the book.
    \item \texttt{year}: The year the book was published.
    \item \texttt{genre}: The genre of the book.
\end{itemize}

Run the following SQL command to create the table (REAL type for float type):

\begin{lstlisting}[language=SQL]
CREATE TABLE books (
    id INTEGER PRIMARY KEY AUTOINCREMENT,
    title TEXT NOT NULL,
    author TEXT NOT NULL,
    year INTEGER,
    genre TEXT
);
\end{lstlisting}

After running the command, you won't see any output, but the table will be created. To verify the table schema, use the `.schema` command:

\begin{lstlisting}[language=bash]
sqlite> .schema books
CREATE TABLE books (
    id INTEGER PRIMARY KEY AUTOINCREMENT,
    title TEXT NOT NULL,
    author TEXT NOT NULL,
    year INTEGER,
    genre TEXT
);
\end{lstlisting}

The `.schema` command displays the exact SQL statement used to create the table, allowing you to confirm its structure.

You use a different command `PRAGMA table\_info` to inspect the structure of a table.
\begin{lstlisting}[language=bash]
sqlite> PRAGMA table_info(books);

\end{lstlisting}

\section*{Step 3: Insert Data into the Table}

Now that we have a table, let's insert some sample data into it.

\begin{lstlisting}[language=SQL]
INSERT INTO books (title, author, year, genre) 
VALUES ('To Kill a Mockingbird', 'Harper Lee', 1960, 'Fiction');

INSERT INTO books (title, author, year, genre) 
VALUES ('1984', 'George Orwell', 1949, 'Dystopian');

INSERT INTO books (title, author, year, genre) 
VALUES ('The Great Gatsby', 'F. Scott Fitzgerald', 1925, 'Classic');

INSERT INTO books (title, author, year, genre) 
VALUES ('Pride and Prejudice', 'Jane Austen', 1813, 'Romance');
\end{lstlisting}

After inserting the data, you won't see any output, but you can verify the data by running a `SELECT` query:

\begin{lstlisting}[language=bash]
sqlite> SELECT * FROM books;
1|To Kill a Mockingbird|Harper Lee|1960|Fiction
2|1984|George Orwell|1949|Dystopian
3|The Great Gatsby|F. Scott Fitzgerald|1925|Classic
4|Pride and Prejudice|Jane Austen|1813|Romance
\end{lstlisting}

\section*{Step 4: Query the Data}

Now that we have some data in the table, let's query it.

\subsection*{1. Select All Books}

To retrieve all the books in the library, run:

\begin{lstlisting}[language=SQL]
SELECT * FROM books;
\end{lstlisting}

The result will look like this:

\begin{lstlisting}[language=bash]
sqlite> SELECT * FROM books;
1|To Kill a Mockingbird|Harper Lee|1960|Fiction
2|1984|George Orwell|1949|Dystopian
3|The Great Gatsby|F. Scott Fitzgerald|1925|Classic
4|Pride and Prejudice|Jane Austen|1813|Romance
\end{lstlisting}

\subsection*{2. Select Books Published After 1900}

To find books published after the year 1900, run:

\begin{lstlisting}[language=SQL]
SELECT * FROM books WHERE year > 1900;
\end{lstlisting}

The result will look like this:

\begin{lstlisting}[language=bash]
sqlite> SELECT * FROM books WHERE year > 1900;
1|To Kill a Mockingbird|Harper Lee|1960|Fiction
2|1984|George Orwell|1949|Dystopian
3|The Great Gatsby|F. Scott Fitzgerald|1925|Classic
\end{lstlisting}

\subsection*{3. Select Books by Genre}

To find all books in the "Fiction" genre, run:

\begin{lstlisting}[language=SQL]
SELECT * FROM books WHERE genre = 'Fiction';
\end{lstlisting}

The result will look like this:

\begin{lstlisting}[language=bash]
sqlite> SELECT * FROM books WHERE genre = 'Fiction';
1|To Kill a Mockingbird|Harper Lee|1960|Fiction
\end{lstlisting}

\section*{Step 5: Update Data}

Let's say we want to update the genre of "1984" from "Dystopian" to "Science Fiction". We can do this with the \texttt{UPDATE} statement:

\begin{lstlisting}[language=SQL]
UPDATE books SET genre = 'Science Fiction' WHERE title = '1984';
\end{lstlisting}

After running the command, you won't see any output, but you can verify the change by querying the data again:

\begin{lstlisting}[language=bash]
sqlite> SELECT * FROM books WHERE title = '1984';
2|1984|George Orwell|1949|Science Fiction
\end{lstlisting}

\section*{Step 6: Delete Data}

If we want to remove a book from the library, we can use the \texttt{DELETE} statement. For example, to delete "Pride and Prejudice":

\begin{lstlisting}[language=SQL]
DELETE FROM books WHERE title = 'Pride and Prejudice';
\end{lstlisting}

After running the command, you won't see any output, but you can verify the deletion by querying the data again:

\begin{lstlisting}[language=bash]
sqlite> SELECT * FROM books;
1|To Kill a Mockingbird|Harper Lee|1960|Fiction
2|1984|George Orwell|1949|Science Fiction
3|The Great Gatsby|F. Scott Fitzgerald|1925|Classic
\end{lstlisting}

\section*{Step 7: Drop the Table (Optional)}

If you want to delete the entire table (including all its data), you can use the \texttt{DROP TABLE} command:

\begin{lstlisting}[language=SQL]
DROP TABLE books;
\end{lstlisting}

After running the command, you won't see any output, but you can verify that the table is gone by listing all tables:

\begin{lstlisting}[language=bash]
sqlite> .tables
\end{lstlisting}

Since the table has been dropped, no tables will be listed.



\section*{Step 8: Other Useful Commands}

You can try other useful commands

\begin{lstlisting}[language=SQL]
.open library.db  
SELECT * FROM books WHERE genre like '%F%';
SELECT * FROM books WHERE lower(genre) like '%f%';
SELECT * FROM books WHERE lower(genre) like '%f%' limit 2;
\end{lstlisting}


\section*{Step 9: Exit SQLite}

Once you're done working with the database, you can exit the SQLite shell by typing:

\begin{lstlisting}[language=bash]
.exit
\end{lstlisting}

The terminal will return to the normal command prompt:

\begin{lstlisting}[language=bash]
sqlite> .exit
user@linux:~$
\end{lstlisting}

\section*{Conclusion}

In this tutorial, we created a simple SQLite database to manage a library system. We learned how to create a table, verify its schema, insert data, query data, update records, and delete records. SQLite is a lightweight, serverless database engine that is perfect for small applications, prototyping, and learning SQL.

Feel free to experiment with more complex queries and operations as you become more comfortable with SQLite!

\end{document}